\chapter*{Global Exercise Bank --- Single Correct MCQs}
\addcontentsline{toc}{chapter}{Global Exercise Bank: Single Correct MCQs}
\label{chap:single_correct}

\begin{tcolorbox}[enhanced,colback=subtlegreen,colframe=cengagegreen,arc=2mm,boxrule=1pt,
    title=\textbf{\large Single Correct Option MCQs (All Chapters)}]
Sources: \textbf{[CNG-TH]} Concept App., \textbf{[GRB]} Ex-I, \textbf{[NA]} Level-1, \textbf{[ESS]}, \textbf{[CNG-DPP]}.
\end{tcolorbox}

\begin{problembox}
\textbf{SC.1} \hfill \textbf{[GRB Ex-I / ESS Q]}

The unit of molar conductance ($\Lambda_m$) is:
\begin{multicols}{2}
\begin{enumerate}[label=(\alph*)]
    \item $\Omega^{-1}$
    \item $\text{S\,cm}^2\,\text{mol}^{-1}$
    \item $\text{S\,cm}^{-1}$
    \item $\Omega\,\text{cm}^2\,\text{mol}^{-1}$
\end{enumerate}
\end{multicols}
\end{problembox}
\begin{solution}
\textbf{(b)} $\text{S\,cm}^2\,\text{mol}^{-1}$ (also written $\Omega^{-1}\,\text{cm}^2\,\text{mol}^{-1}$ or $\text{S\,m}^2\,\text{mol}^{-1}$ in SI). $\Lambda_m = \kappa/c$ where $\kappa$ (S\,cm$^{-1}$), $c$ (mol\,cm$^{-3}$).
\end{solution}

\begin{problembox}
\textbf{SC.2} \hfill \textbf{[CNG-TH Concept App. / GRB]}

For a weak electrolyte at infinite dilution, which of Kohlrausch's law gives $\Lambda_m^\circ$?
\begin{enumerate}[label=(\alph*)]
    \item By extrapolating the Kohlrausch plot ($\Lambda_m$ vs $\sqrt{c}$) to $c=0$
    \item By adding $\lambda^\circ$ of constituent ions algebraically
    \item Both (a) and (b)
    \item (a) only for strong electrolytes; (b) only for weak electrolytes
\end{enumerate}
\end{problembox}
\begin{solution}
\textbf{(d)} For strong electrolytes, the $\Lambda_m$ vs $\sqrt{c}$ plot is linear and can be extrapolated to $c\to0$ (method a). For weak electrolytes (like acetic acid), the plot shows a sharp, non-linear upswing at low concentrations and cannot be extrapolated reliably; instead, Kohlrausch's law of independent migration (method b) must be used to calculate $\Lambda_m^\circ$ indirectly from strong electrolyte data.
\end{solution}

\begin{problembox}
\textbf{SC.3} \hfill \textbf{[NA Level-1 Q6]}

The molar conductance of $0.1\,\text{M}$ \ce{CH3COOH} is $5.2\,\text{S\,cm}^2\,\text{mol}^{-1}$ and $\Lambda_m^\circ(\ce{CH3COOH}) = 390.7\,\text{S\,cm}^2\,\text{mol}^{-1}$. The degree of dissociation $\alpha$ is:
\begin{multicols}{4}
\begin{enumerate}[label=(\alph*)]
    \item $0.0133$
    \item $0.133$
    \item $0.75$
    \item $0.33$
\end{enumerate}
\end{multicols}
\end{problembox}
\begin{solution}
\textbf{(a)} $\alpha = \Lambda_m/\Lambda_m^\circ = 5.2/390.7 = 0.0133$
\end{solution}

\begin{problembox}
\textbf{SC.4} \hfill \textbf{[GRB Ex-I / CNG-TH]}

In the standard hydrogen electrode (SHE), the standard conditions are:
\begin{enumerate}[label=(\alph*)]
    \item $[\ce{H+}] = 1\,\text{mol\,L}^{-1}$, $P_{\ce{H2}} = 1\,\text{atm}$, $T = 298\,\text{K}$
    \item $[\ce{H+}] = 1\,\text{mol\,L}^{-1}$, $P_{\ce{H2}} = 1\,\text{bar}$, $T = 298\,\text{K}$
    \item $a_{\ce{H+}} = 1$, $f_{\ce{H2}} = 1\,\text{bar}$, $T = 298\,\text{K}$
    \item $[\ce{H+}] = 7\,\text{M}$, $P_{\ce{H2}} = 1\,\text{bar}$, $T = 298\,\text{K}$
\end{enumerate}
\end{problembox}
\begin{solution}
\textbf{(c)} The strictly correct thermodynamic definition uses \textit{activities}: $a_{\ce{H+}} = 1$ (unit activity of H$^+$) and $f_{\ce{H2}} = 1\,\text{bar}$ (unit fugacity of H$_2$). In practice, $1\,\text{mol\,L}^{-1}$ \ce{HCl} approximates unit activity, but for a rigorous JEE/Olympiad answer, (c) is most correct.
\end{solution}

\begin{problembox}
\textbf{SC.5} \hfill \textbf{[CNG-DPP DPP-2 / GRB Q]}

Which of the following couples has the \textbf{highest} tendency to act as a reducing agent?
\begin{multicols}{2}
\begin{enumerate}[label=(\alph*)]
    \item $\ce{Li+/Li}$ ($E^\circ = -3.05\,\text{V}$)
    \item $\ce{Zn^2+/Zn}$ ($E^\circ = -0.76\,\text{V}$)
    \item $\ce{Cu^2+/Cu}$ ($E^\circ = +0.34\,\text{V}$)
    \item $\ce{Ag+/Ag}$ ($E^\circ = +0.80\,\text{V}$)
\end{enumerate}
\end{multicols}
\end{problembox}
\begin{solution}
\textbf{(a)} The most negative $E^\circ$ means the metal is most easily oxidized (strongest reducing agent). Li$^+$/Li has $E^\circ = -3.05\,\text{V}$: Li is the strongest reducing agent in the electrochemical series.
\end{solution}

\begin{problembox}
\textbf{SC.6} \hfill \textbf{[ESS §24 / GRB Q]}

The transference number of $\ce{K+}$ in \ce{KCl} at infinite dilution is approximately $0.49$. This means:
\begin{enumerate}[label=(\alph*)]
    \item $49\%$ of the current is carried by \ce{K+} ions
    \item $49\%$ of the ions are \ce{K+}
    \item \ce{K+} moves $49\%$ faster than \ce{Cl-}
    \item $49\%$ of \ce{KCl} is dissociated
\end{enumerate}
\end{problembox}
\begin{solution}
\textbf{(a)} The transference number $t_+ = 0.49$ means $49\%$ of the electric current is transported by \ce{K+} ions; the remaining $51\%$ is transported by \ce{Cl-} ions ($t_- = 0.51$). Since $t_+ \approx t_-$, both ions have nearly equal mobilities --- hence KCl (and KNO$_3$, NH$_4$NO$_3$) are used in salt bridges.
\end{solution}

\begin{problembox}
\textbf{SC.7} \hfill \textbf{[NA Level-1 Q12 / GRB]}

The $\text{EMF}$ of the cell $\ce{Zn | ZnSO4(0.01 M) || CuSO4(1 M) | Cu}$ at $298\,\text{K}$ is: [$E^\circ(\ce{Cu^2+/Cu})=+0.34\,\text{V}$, $E^\circ(\ce{Zn^2+/Zn})=-0.76\,\text{V}$]
\begin{multicols}{2}
\begin{enumerate}[label=(\alph*)]
    \item $1.10\,\text{V}$
    \item $1.16\,\text{V}$
    \item $1.04\,\text{V}$
    \item $0.98\,\text{V}$
\end{enumerate}
\end{multicols}
\end{problembox}
\begin{solution}
\textbf{(b)}\\
$E^\circ_{\text{cell}} = 0.34 - (-0.76) = 1.10\,\text{V}$\\
Nernst ($n=2$): $E = 1.10 - \frac{0.0591}{2}\log\frac{0.01}{1} = 1.10 - \frac{0.0591}{2}\times(-2) = 1.10 + 0.0591 = \mathbf{1.16\,\text{V}}$
\end{solution}

\begin{problembox}
\textbf{SC.8} \hfill \textbf{[CNG-TH Concept App / NA Q14]}

Given $K_{sp}(\ce{AgCl}) = 1.8\times10^{-10}$ at $298\,\text{K}$ and $E^\circ(\ce{Ag+/Ag}) = +0.80\,\text{V}$, the standard reduction potential of the $\ce{Ag | AgCl(s) | Cl-}$ electrode is:
\begin{multicols}{2}
\begin{enumerate}[label=(\alph*)]
    \item $+0.22\,\text{V}$
    \item $+0.50\,\text{V}$
    \item $-0.22\,\text{V}$
    \item $+0.80\,\text{V}$
\end{enumerate}
\end{multicols}
\end{problembox}
\begin{solution}
\textbf{(a)}\\
$E^\circ(\ce{AgCl/Ag}) = E^\circ(\ce{Ag+/Ag}) + 0.0591\log K_{sp}$\\
$= 0.80 + 0.0591\log(1.8\times10^{-10}) = 0.80 + 0.0591\times(-9.745) = 0.80 - 0.576 = \mathbf{+0.224\,\text{V} \approx +0.22\,\text{V}}$
\end{solution}

\begin{problembox}
\textbf{SC.9} \hfill \textbf{[GRB / CNG-DPP DPP-3]}

In a Daniell cell, the salt bridge contains saturated \ce{KCl}. Its primary function is to:
\begin{enumerate}[label=(\alph*)]
    \item Provide \ce{K+} and \ce{Cl-} as reactants in the cell reaction
    \item Prevent mixing of the two electrolytes
    \item Maintain electrical neutrality in both half-cells and eliminate liquid junction potential
    \item Increase the cell potential by providing a third electrolyte
\end{enumerate}
\end{problembox}
\begin{solution}
\textbf{(c)} The salt bridge serves two roles: (i) maintains electrical neutrality by allowing ion migration (\ce{K+} into the cathode solution, \ce{Cl-} into the anode solution, as ions are consumed/produced), and (ii) minimizes the liquid junction potential because $t_+ \approx t_-$ for \ce{KCl} (both ions have nearly equal transference numbers), so the junction potential nearly cancels.
\end{solution}

\begin{problembox}
\textbf{SC.10} \hfill \textbf{[NA Level-1 Q20]}

The equilibrium constant for the cell reaction of the Daniel cell (\ce{Zn | Zn^2+ || Cu^2+ | Cu}) at $298\,\text{K}$ is: [$E^\circ_{\text{cell}} = 1.10\,\text{V}$, $n=2$]
\begin{multicols}{2}
\begin{enumerate}[label=(\alph*)]
    \item $10^{37.2}$
    \item $10^{3.72}$
    \item $e^{85.6}$
    \item $10^{18.6}$
\end{enumerate}
\end{multicols}
\end{problembox}
\begin{solution}
\textbf{(a)}\\
$\log K = \frac{n\,E^\circ}{0.0591} = \frac{2\times1.10}{0.0591} = \frac{2.20}{0.0591} = 37.23$\\
$K = 10^{37.23} \approx \mathbf{10^{37.2}}$
\end{solution}
